\section{Source audit, proof checks, and reproducibility}
\label{sec:audit}

\subsection{The five incoming reports}

The source audit distinguishes the following archives at the pinned
commit.  The short names in this section are only reading aids; the
machine-readable record preserves every full filename and internal path.
\begin{description}[style=nextline,leftmargin=13mm]
\item[RC] \nolinkurl{ProveIt_Resonance_Correlations_2026-10-10.zip}.
Existing Gamma and Hurwitz-jet correlations supply background.
The present triple formula answers the coordinate part of item G3
in \nolinkurl{sections/07_questions.tex}; minimal character spaces
and arithmetic independence remain separate.
\item[RJ] \nolinkurl{ProveIt_Shifted_Hurwitz_Jets_2026-10-10.zip}.
Its bilinear kernel and unit-scale contact terms are retained.
The questions on three shifted factors and on dilation, traces, and
subtraction scales in \nolinkurl{sections.tex} are answered here under
the stated separation and positive-integer hypotheses.
\item[SC] \nolinkurl{ProveIt_Stieltjes_Convolution_2026-10-10.zip}.
Its periodic convolution algebra and one-factor derivative anomalies
are established input.  The affine-frequency future-work sentence
requires the qualification described below.
\item[CL] \nolinkurl{ProveIt_Stieltjes_Correlation_Closure_2026-10-10.zip}.
Its all-index bilinear linear closure and missing-order rule are
rederived compactly in Theorem~\ref{base:theorem}.  The trilinear and
unequal-integer-dilation questions in its final section are answered
for separated singularities.
\item[CC] \nolinkurl{Stieltjes_Convolution_Calculus_2026-10-10.zip}.
Its derivative and primitive ladders overlap the other reports.
The first further-research item, concerning a triple generating
function and subtraction rule, is answered here; the period and
coordinate question is answered for affine integer rescaling.
\end{description}

The exact anchors and Git blob identifiers appear in
\nolinkurl{integration/source_audit.json}.  Overlap is recorded as
established background, not counted again as a present result.

\subsection{Proposed source corrections}

The only definite correction proposed to the reviewed repository text
is to SC's future-work discussion, \nolinkurl{article.tex},
lines 1118--1121.  Two nonzero integer affine frequencies leave one free
integer in $pn+qm=0$.  Thus a different integer frequency alone does not
force a higher-dimensional summation index.  The replacement text in
\nolinkurl{integration/proposed_corrections.tex} states the precise
bilinear theorem and reserves the higher-dimensional issue for three
or more independent factors.  This correction does not invalidate
SC's convolution identities.

The external corrections to Choi's printed p.~6 are mathematically
different: the displayed identities (41)--(46), with the stated
definitions, have twice the correct right side, and two harmonic
orders are misprinted.  Theorem~\ref{gs:choi-master} gives corrected
formulas and a proof for every order.  The correction record includes
the DOI, page, and equation numbers.  The publisher PDF was visually
checked.  Neither the absence of a located erratum nor numerical
agreement is used to infer historical priority.

The independent proof audit of this article checked the Fourier phase
convention, the six disjoint frequency cones, the pole residues, the
Gamma and harmonic factorials, and the mixed parameter domains.  It
also sharpened the uniform local estimate used for primitive
correlations and the Cauchy argument at simultaneous spectral poles.
These improvements are incorporated into the proofs as delivered.

% Optional manuscript fragment. All formulas below are existing conjectures
% in the pinned source; none is asserted to be proved by the present report.
\subsection{The surviving named Gaussian conjectures}
\label{sec:surviving-gaussian-conjectures}

The manuscript's names $S_6$ and $S_8$ refer to specific alternating
harmonic sums, and should not be confused with the uncolored sums of our
Gamma-ratio theorem.  With the largest summation index first in the
multiple-polylogarithm convention, their definitions are
\[
 S_p=\sum_{n\geq0}\frac{(-1)^nH_n}{(2n+1)^p},\qquad
 g_{a,b}=\operatorname{Im}\operatorname{Li}_{a,b}(i,1)
 =\sum_{n\geq0}\frac{(-1)^nH_{2n}^{(b)}}{(2n+1)^a},
\]
where $H_m^{(b)}=\sum_{j=1}^{m}j^{-b}$ and $H_m=H_m^{(1)}$.
Write $G=\beta(2)$ for Catalan's constant.  The existing mixed-endpoint
identity gives
\[
 S_p=g_{p,1}+\operatorname{Im}\operatorname{Li}_{p,1}(i,-1).
\]
In particular, $S_p$ is not just $g_{p,1}$.

The frozen weight-seven conjecture currently reads
\begin{align}
 S_6\stackrel{?}{={}}&\frac{722}{527}g_{6,1}
 +\frac{40}{527}g_{4,3}-\frac{128}{155}g_{2,5}
 +\frac{15191}{28569600}\pi^7\notag\\
 &-\frac3{124}G\zeta(5)
 -\frac{2373}{10540}\beta(4)\zeta(3)-2\beta(6)\log2.
 \label{eq:source-current-S6}
\end{align}
The source is \nolinkurl{chapters/04-cyclotomic-quotients.tex},
lines 338--349.  Its label is \texttt{cycloquot:conj:S6}.
The alternative formula involving $g_{5,2}$ is already proved equivalent
to this one by two convergent shuffle identities in the same file,
lines 295--336.  That coordinate change does not prove either conjectured
vanishing.  Moreover, the source's separating functional is nonzero on
the frozen relation, so a proof cannot use only the restricted
single-product presentation considered there.

The current weight-nine conjecture is
\begin{align}
 S_8\stackrel{?}{={}}&\frac{52334}{37973}g_{8,1}
 +\frac{2824}{37973}g_{6,3}
 -\frac{6832}{189865}g_{4,5}
 -\frac{139760}{265811}g_{2,7}\notag\\
 &+\frac{998183\pi^9}{17283022848}
 -\frac{78615}{19442176}G\zeta(7)
 -\frac{201021}{4252976}\beta(4)\zeta(5)\notag\\
 &-\frac{290505}{1063244}\beta(6)\zeta(3)-2\beta(8)\log2.
 \label{eq:source-current-S8}
\end{align}
Its exact source is \nolinkurl{chapters/04-S8-candidate.tex},
lines 7--17, label \texttt{s8new:conj:S8}.
The same file, lines 40--50, records exact rational proximity bounds
below $10^{-355}$ for the normalized residuals of both current
conjectures.  Both enclosures contain zero.  The equations remain
conjectural in the pinned manuscript and in the present report.

This $S_8$ candidate must be distinguished from the earlier frozen vector
in the basket containing $g_{7,2}$ and $g_{5,4}$.  That different vector
has a strictly negative certified residual between
$-1.3274179020580459\times10^{-86}$ and
$-1.3274179020580457\times10^{-86}$; see
\nolinkurl{chapters/10-discovery.tex}, lines 259--301,
label \texttt{research:prop:S8-rejected}.
The rejection concerns precisely that vector and gives no arithmetic
independence result.

Our uncolored harmonic theorem does not acquire the factor $(-1)^n$
merely by setting its Hurwitz parameter to $1/2$.  Its direct application
therefore cannot prove either displayed Gaussian identity.  The
trilinear theorem does provide root-of-unity colored Tornheim
coordinates at rational shifts, with multiple-polylogarithm reductions
at positive integer inner orders.  A coordinate representation alone
does not annihilate either frozen residual, and the partial-fraction
identities for discrete inner orders cannot be differentiated as though
they held with a continuously varying summation range.  The remaining
concrete problem is to find and prove the missing mixed-endpoint
functional relation at weights seven and nine, and then reduce it
exactly to the specified frozen basket.


\subsection{Independent numerical diagnostics}

The programs evaluate two mathematically different sides whenever
practical.  They are numerical diagnostics, not interval certificates;
the proofs in the preceding sections establish the identities.
Working precision is larger than the number of digits reported.

\begin{center}
\small
\begin{tabularx}{\textwidth}{@{}>{\raggedright\arraybackslash}p{30mm}X >{\raggedright\arraybackslash}p{29mm}@{}}
\toprule
Component & Independent comparison & Largest observed residual\\
\midrule
Integer dilations & Local digamma subtraction against the scale formula;
four frequency pairs at 60 digits & $7.4\times10^{-58}$\\
Ordinary log-Gamma & Direct split quadrature against polylogarithmic
order finite differences, step $10^{-7}$ & $1.9\times10^{-31}$\\
Triple products & Local three-digamma subtraction against a separate
Mellin and polylogarithm-order evaluation, at 42 working digits
& $8.7\times10^{-42}$\\
Shifted harmonic sums & 37 Mellin and local-subtraction comparisons
against explicit polygamma formulas, at 70 digits
& $1.82\times10^{-56}$\\
Gauss mixed jets & 48 direct sums with Bernoulli-polynomial tails
against mixed Gamma coefficients, at 75 digits
& $5.14\times10^{-58}$\\
Negative balance & Direct central-binomial bracket with an independent
tail, at 75 digits & $1.97\times10^{-58}$\\
\bottomrule
\end{tabularx}
\end{center}

The reported maxima are rounded upward from the stored results.  For
the Gauss checks the change when the tail parameters are refined from
$(N,K)=(70,24)$ to $(110,32)$ is below $6.00\times10^{-42}$.
This refinement is useful evidence about numerical stability and is
not itself a rigorous error estimate.  The log-Gamma finite-difference
comparison has deliberately lower accuracy than the digamma finite
part; its error is dominated by the derivative approximation.

The dilation check computes the unit-frequency base integral by its own
local subtraction.  Consequently it checks the new frequency and scale
formula without assuming the closed evaluation of $J_{00}$.  The
triple check additionally compares the six-cone formula at
$\boldsymbol\lambda=(1,1,1)$ and $(2,1,3)$, with rational shifts
$0,1/5,7/10$, against exact piecewise Bernoulli-polynomial integration.
The exact values are respectively $-1/1000$ and
$-124409/7560000000$.

For the equal-thirds digamma example, the two numerical routes give
\begin{equation}\label{eq:numeric-triple}
 \FP\int_0^1\psi(x)\psi(\{x+\tfrac13\})
                     \psi(\{x+\tfrac23\})\,dx
 =27.805592568749290182019060698117\ldots.
\end{equation}
This decimal is a numerical evaluation of the exact Tornheim formula
\eqref{eq:triple-third}, not a replacement for it.

\subsection{Exact arithmetic and reproducible artifacts}

The exact coefficient generator in
\nolinkurl{code/exact_coefficients.py} uses SymPy arithmetic to produce
Gamma coefficients, rational-kernel reductions, local contact
polynomials, bilinear Stieltjes coefficients, and the corrected
two-parameter harmonic sums.  Its inputs and outputs are exact; no
integer-relation search is used.  The formal Stieltjes symbols in a
coefficient table are named atoms, with no assumed arithmetic
independence.  A separate exact six-cone check is explained in
Appendix~\ref{app:exact}.

The included scripts and JSON results make each numerical comparison
reproducible.  \nolinkurl{README.md} gives the build and check commands,
\nolinkurl{requirements.txt} records the Python packages, and
\nolinkurl{MANIFEST.sha256} records checksums of the final files.
\nolinkurl{integration/INTEGRATION_NOTES.md} maps theorem labels to
source questions and lists the necessary finite-part conventions.
These files are a research supplement.  No Lean, Isabelle, or other
proof-assistant certification is claimed.
